\section{The source layers}
\label{sec:sources}

Every derivation in this program can be traced to bytes. This section
says how, because the tracing is not uniform: there are three source
systems, they certify different things, and reading a statement out of
the wrong one is the most common way to get this corpus wrong.

An agent that internalizes only one thing from this document should
internalize the table below.

\begin{center}
\begin{tabular}{@{}p{0.19\linewidth} p{0.34\linewidth} p{0.38\linewidth}@{}}
\toprule
Layer & What it holds & What it certifies \\
\midrule
\textbf{Theory graph} &
\DerivStatements{} derivation statements in \CiteDocs{} documents;
\ResTotal{} results; \LeanTotal{} theorems; \RouteTotal{} routes;
\GraphEdges{} edges in \GraphEdgeKinds{} relations &
That a statement was \emph{recorded} with a status, hypotheses, a
source locator, and a digest of the bytes it was read from. It does
not certify the mathematics; the tier does that separately. \\
\addlinespace
\textbf{Published work} &
\LitPapers{} papers in the evidence map, \LitVerified{} edges verified;
\LibraryPapers{} acquired papers in the reading library;
\ExtractedPages{} extracted pages across \ExtractedWorks{} works &
That an external result exists and, for a verified edge, that it was
read against a named claim here. Obtained is not read; read is not
checked. \\
\addlinespace
\textbf{Local archive} &
Four years of manuscripts, runs, notes and Lean, with
\AliasCount{} aliased documents, \NoteArchives{} note archives and
\NoteReviews{} recorded reviews &
That the bytes are preserved and pinned. Nothing about whether the
content is current, or correct. \\
\bottomrule
\end{tabular}
\end{center}

\subsection{Layer 1: the theory graph}

The graph is the index of record. Its node identifiers are prefixed,
and the prefix tells you what kind of object you are holding and
therefore what it can support.

\begin{description}[leftmargin=1.1em,style=nextline]
\item[\texttt{CITE:}] A source document, pinned by path and SHA-256.
  \CiteDocs{} of them. A \texttt{CITE:} record is bytes, not a claim.
\item[\texttt{DERIV:}] A statement located \emph{inside} a
  \texttt{CITE:} document, by heading and line range.
  \DerivStatements{} of them, every one with a locator. This is the
  finest granularity the corpus tracks, and it is the right unit to
  work at.
\item[\texttt{RESULT:}] A first-class analytic claim, independent of any
  one document. \ResTotal{} of them, each with status, evidence,
  hypotheses, scope, \texttt{depends\_on} and \texttt{bears\_on}.
\item[\texttt{LEAN:}] A registered theorem. \LeanTotal{} of them.
  \texttt{formalizes} links it to a catalogue record;
  \texttt{promotes} --- which only 49 carry --- means it proves an
  entire named check and therefore lifts that check's tier.
\item[\texttt{CHK:}] A registered invariant check, carrying its tier and
  the numbers it establishes. A \texttt{FINDING:} check is one that
  asserts a \emph{discrepancy} rather than an agreement.
\item[\texttt{RUN:}] A pinned computational run, with its own manifest.
\item[\texttt{ADR:}] A numbered decision record: what was decided, on
  what evidence, and what it retracts or amends.
\item[\texttt{CONST:}, \texttt{SYM:}] A named constant and its symbol
  registry entry --- canonical name, corpus spellings, code names, and
  exact values. Appendix~\ref{app:symbols}.
\item[\texttt{G}$n$, \texttt{C}$n$, \texttt{U}$n$, \texttt{R}$n$] Gaps,
  contradictions, unifying candidates and registry claims.
\end{description}

The edges matter as much as the nodes. \GraphEdges{} of them across
\GraphEdgeKinds{} relations, and three deserve separate mention because
they are routinely confused. \texttt{depends\_on} is a mathematical
input: if it falls, the dependent falls. \texttt{supported\_by} is
narrower --- a control that bears on part of a statement, and a result
with \texttt{supported\_by: []} has no registered machine support at
all, which is a fact worth checking before quoting a tier.
\texttt{bears\_on} is looser still: a topical connection, and it
establishes nothing.

\subsection{Layer 2: published work}

There are two published-work systems and they answer different
questions. Using the wrong one produces either a false novelty claim or
a false authority claim.

\emph{The evidence map} (\srcpath{literature/index.yaml}) is small,
curated and adversarial. \LitPapers{} papers, \LitEdges{} recorded
relations, \LitVerified{} of them verified against a named claim of
this program. An entry earns its place by naming a specific claim.
Query it with \texttt{workhouse lit -{}-for G19}. This is the only
layer from which an agreement or a contradiction may be quoted.

\emph{The reading library}
(\srcpath{literature/library/paper\_manifest.csv} and
\srcpath{volume2/}) is large and inclusive: \LibraryPapers{} acquired
papers --- \LibraryVolOne{} and \LibraryVolTwo{} --- each with tier,
category, journal reference, DOI, arXiv identifier, INSPIRE citation
count, local file, page count, SHA-256, abstract, and a curated
one-line \texttt{workhouse\_connection}. Appendix~\ref{app:library}
reproduces it. Its purpose is coverage: to answer ``has anyone done
this?'' before an agent spends a week. Its connection field is a
reading note and carries no evidential weight.

The gap between the two is the useful measurement. \LitObtained{}
papers have been read against a claim; \LibraryPapers{} have been
acquired. Most of the library has never been checked against anything
here, and a novelty claim asserted against the library is therefore
weaker than one asserted against the evidence map --- which is itself
weaker than priority, for the reasons in \S\ref{sec:external}.

\subsection{Layer 2b: the page-level extraction index}

This is the piece an agent is most likely to miss, and it is the one
that saves the most time.

\ExtractedPages{} pages of local PDFs have been extracted and indexed
at \srcpath{literature/yangmills/extraction/pages.jsonl}, covering
\ExtractedWorks{} works. Each record carries the paper identifier, the
PDF page number, the SHA-256 of the source \emph{and} of the extracted
text, the path of the reading copy, and the line positions of every hit
against \ExtractedTopics{} topic keys. Appendix~\ref{app:extraction}
tabulates it.

The workflow it supports:

\begin{enumerate}
  \item Search \srcpath{pages.jsonl} for a paper identifier and a topic
    key (\texttt{ground\_state}, \texttt{spectral\_gap},
    \texttt{curvature}, \texttt{reconstruction},
    \texttt{uniform\_limits}, \texttt{constructive},
    \texttt{gauge\_geometry}).
  \item Follow the page and line locator into the reading copy under
    \srcpath{literature/inbox/<ID>/extracted/}.
  \item Verify any equation you intend to use against the pinned PDF.
\end{enumerate}

Three cautions the index states about itself. A topic hit may be a
citation rather than a result, or an extraction error. Candidate
theorem labels are recorded on many pages, but a label can sit inside a
citation or carry an OCR fault --- the index is a search aid and
asserts nothing. And five pages extracted fewer than eighty characters
and are explicitly flagged as unreliable; their contents must not be
inferred from the text index. Exact equations from legacy scans must be
checked in the PDF.

\subsection{Layer 3: the local archive}

The archive is preserved rather than curated, and its governing
instruction is that nothing is erased. Failed attempts, superseded
drafts, duplicates and uncommitted work are all kept, because deleting
a refuted claim destroys the evidence that it was tried.

Consequently the field that matters when reading an archive document is
its \emph{standing}. \AliasCount{} documents carry a registered alias
and a standing, tabulated in Appendix~\ref{app:documents}. A superseded
document is not marked in its own text and reads exactly like a current
one; the register is the only thing that distinguishes them. Anything
under \srcpath{theory/superseded/} is retained for audit and must never
be read as current.

The maintainer's own research notes enter through a separate
discipline: \NoteArchives{} declared archives with \NoteReviews{}
recorded reviews, each carrying a content digest, a verdict from a
closed vocabulary, and a mandatory reason. An \texttt{import} is
byte-verified against its digest; an \texttt{extract} must name the
claims it entered through; a \texttt{set-aside} is a recorded judgement
and never a deletion. No verdict promotes anything past the default
machine tier.

\subsection{Finding things: the corpus joins on values, not concepts}

The single most important retrieval fact about this corpus:
\emph{its join keys are exact rationals and spellings, not concepts.}
No semantic query retrieves $109151/249696$, and the same quantity
appears in a manuscript, in code, and in a claim record under three
different names.

\texttt{workhouse search} exists for this. It matches by \emph{value}
rather than spelling --- searching $-10/96$ finds $-5/48$ --- across
the prose index, the code and certificate index, the claim catalogue,
and the curated alias table of Appendix~\ref{app:symbols}. It also
carries two warnings no text search can: names the register
\emph{forbids} because using them regenerates a dissolved
contradiction, and names coined inside this program that its own source
corpus never uses.

Two indexes sit underneath it, and the second exists because the first
could not see the sealed core: values such as $5/48$, $5/12$, $5/612$,
$11/306$ and $7/102$ can be absent from the prose index while their
computations are present in code.

\subsection{The queries worth knowing}

\begin{verbatim}
workhouse why <id>          everything recorded about one claim:
                            both sides of a dispute, every check and
                            its verdict, routes tried and which died
workhouse search '5/612'    search by exact VALUE, not by name
workhouse branches C2       both sides of a dispute, side by side
workhouse derive C2 G3      the evidence chain, registered edges only
workhouse lit --for G19     published work bearing on a claim
workhouse notes --queue     the next unreviewed notes, by signal
workhouse triage <path>     survey an unpinned collection
workhouse ask "<question>"  candidate finder over prose; ends in `why`
workhouse brief <gap> --json  target snapshot, with freshness
\end{verbatim}

\texttt{why} is the front door. It resolves a claim identifier, a
constant, a gap, an ADR, or a corpus file path, and it prints the
routes with their states --- which is how a dead route announces itself
before an agent re-runs it.

\subsection{What each layer cannot do}

\begin{itemize}
  \item The graph cannot tell you a statement is \emph{true}. It tells
    you what was recorded, by whom, from which bytes, under which
    hypotheses. Status, evidence and machine tier are three
    independent fields precisely so that this stays visible.
  \item The library cannot establish novelty. It establishes that a
    paper was acquired.
  \item The extraction index cannot establish what a paper says. It
    tells you which page to read.
  \item The archive cannot tell you what is current. Standing does.
  \item A digest cannot tell you an anchor resolves. Check both: a
    correct SHA-256 over a fabricated locator has occurred here.
\end{itemize}

\subsection{Operational rules live elsewhere, deliberately}

How to add an invariant, register a Lean theorem, record a route, order
the regeneration passes, and land work: those are specified in
\srcpath{CLAUDE.md} (the working agreement), \srcpath{AGENTS.md} (the
research posture) and \srcpath{CONTRIBUTING.md}. They are not restated
here.

That is a decision, not an omission. Those files are binding and they
change; a copy in a document that is rebuilt on its own schedule would
drift, and an agent following a stale copy of a rule is worse off than
one following no copy. This document carries the \emph{mathematical}
working knowledge --- which sources support what, how the derivations
connect, and where the traps are. For procedure, read the source of
record.

Two of its rules are worth naming here anyway, because they explain the
shape of everything above. A claim is established by its proof or
computation under explicit hypotheses, and its artefact must not claim
more than that argument establishes. And a disputed value closes by
derivation or not at all: both sides stay recorded, and no code path
may pick one, average them, or prefer whichever is an exact rational
because exactness reads as authority.
